\section{Appendix}

\subsection{Validation summary}

The country-screening release should be treated as defensible only if the following checks pass:

\begin{qualitycheck}
\begin{enumerate}
  \item Required columns exist in the country-month and component marts.
  \item Scoring rows have no null values in required fields.
  \item Country coverage matches across marts.
  \item Country-year coverage is complete for years used in scoring.
  \item Duplicate country-month keys are not present after filtering.
  \item Score components remain within expected normalized bounds.
  \item Sensitivity analysis confirms whether the top-five shortlist is stable enough for a screening decision.
\end{enumerate}
\end{qualitycheck}

\subsection{Sensitivity interpretation}

The score weights are judgmental, so the business recommendation should focus on shortlist stability rather than precise rank order. If a country enters the top five only under a growth-led scenario, it should be treated as a watchlist case rather than an immediate Phase 2 priority.

\subsection{Limitations}

\begin{limitation}
Country-level aggregation hides destination-level heterogeneity. A country may rank highly because of one or two cities, because of a broad regional pattern, or because multiple destination types are averaged together. These cases require different commercial responses and cannot be distinguished at country level.
\end{limitation}

\begin{limitation}
The model does not yet include campaign cost, supply availability, regulation, price elasticity, competitive intensity, or historical campaign response. These variables must be introduced before final investment decisions.
\end{limitation}

\subsection{Reproducibility commands}

\begin{terminal}
uv sync
uv run python main.py
uv run marimo edit notebooks/02_country_screening.py
quarto render reports/technical/country_screening_methodology.qmd
cd reports/business && make pdf
\end{terminal}
